\section{Optical Test of Quantum Mechanics}
Today we have the capability to perform experiments on single atoms, molecules, or ions, contradicting what Schrödinger thought would be possible. Fundamental tests are tests that distinguish between quantum mechanical theory and hidden-variable theories.

\subsection{Photon Sources: Spontaneous Parametric Down-Conversion}
We consider the non-degenerate case for parametric down-conversion with the interaction Hamiltonian $\hamiltonian_I \sim \chi^{(2)}\a_p \ad_s\ad_i + \text{h.c.}$ with $p,s,i$ being the pump, signal and idler. Imagining the simplest case $\ket{1}_p\ket{0}_s\ket{0}_i \rightarrow \ket{0}_p\ket{1}_s\ket{1}_i$, where the two photons are created almost instantaneously in the signal and idler, and which is called spontaneous since the signal and idler originally are in the vacuum state. We must respect conservation and thus, $\omega_p = \omega_s + \omega_i$ and inside the crystal we must have $\vec{k}_p \approx \vec{k}_s + \vec{k}_i$, where the approximation stems from the uncertainty in crystal length. These two conditions are called the phase-matching conditions and only non-centrosymmetric crystals have a non-vanishing $\chi^{(2)}$. 

There are two types of spontaneous parametric down-conversion. In type I the signal and idler have the same polarization and which is orthogonal to the pump. The outgoing light will be two concentric cones, one for the signal and one for the idler. A single down-conversion will come out on opposite sides of the two cones. That is, if the signal comes out at $\pi/2$ (measured from the $x$-axis) the idler comes out at $3\pi/2$. In type II the signal and idler are orthogonal, which creates two off-centre cones, the ordinary and extraordinary ray, with an intersection in two points. Photons emerging at these two points will be polarization entangled. \question{I don't really understand why neither cone is not concentric on the pump axis. It should have something to do with the phase-matching I think?} Type II down-conversion may be used to create the bell states $\ket{\Psi^\pm} = (\ket{H}_1 \ket{V}_2 \pm \ket{V}_1 \ket{H}_2)/\sqrt{2}$ and $\ket{\Phi^\pm} = (\ket{H}_1\ket{H}_2 \pm \ket{V}_1\ket{V}_2)/\sqrt{2}$.

\subsection{The Hong-Ou-Mandel Interferometer}
Two single-photon states incident on a beam splitter will experience bunching effects, and detectors will never measure coincidences. An interpretation of this is that the photons are incident on the beam splitter simultaneously. \question{I do not understand this interpretation.} Hong-Ou-Mandel tested this effect by using type I down-conversion. By creating two photons using down-conversion and then aiming them on mirrors to a beam splitter and moving the beam splitter slightly one could measure the time-difference of arrival. Hong-Ou-Mandel (HOM) did this experiment with great success and manages to confirm theory. They did not reach 0 coincidences with the beam splitter at equal distances, which is contributed to an alignment error of the beams.

\subsection{The Quantum Eraser}
The photons in the HOM experiment are being created by the same source of the type I down-conversion and thus have the same polarization, and being indistinguishable is the reason for the interference. Put another way, the inability to detect which path the photon has taken causes interference. Placing a rotator at the idler beam and rotating it to vertical polarization will completely destroy the interference, causing classical statistics to dictate the count rate ($P_\text{coincidence} = 0.5$). Even if the polarization is not measured in the detector, the possibility of detecting polarization, and thus which path the photon took, removes the interference.

Introducing linear polarizers at the detectors with angles $\theta_1,\theta_2$, we get the probability to detect coincidences $P_\text{coincidence} = \sin^2(\theta_2 - \theta_1)/4$. That is, we can choose the angles such that $P_\text{coincidence} = 0$, regaining the bunching effect and the interference. We have erased the path-information. This work by the polarizers projecting the state onto two orthogonal states, thus removing the possibility of double counts. We note that the process of erasing the path-information take place outside the HOM-interferometer and after the photons have crossed paths. That is, the photons do not interact in any way again after the polarizers.

\stepcounter{subsection}\subsection{Superluminal Tunnelling of Photons}
Imagining once again a balanced HOM-interferometer, but now with a dielectric mirror, acting as a tunnelling barrier at the signal beam. Interestingly, by the insertion of this barrier to regain a balance in coincidence, we must lengthen the signal path, not the idler path. The important distinction being that the tunnelling through the mirror happens superluminally. The common explanation for this phenomenon is that the transmitted Gaussian wave packet is created by pulse reshaping of the front of the Gaussian wave packet before interacting with the barrier. That is, the front of the pulse is being reshaped to new Gaussian pulse and transmitted before the peak of the original pulse have hit the barrier. This can not be used for faster-than-light signalling. \question{Is this the same as fast light that can be seen in materials with negative refractive indices. For example Rare-Earth ion-crystals.}

\subsection{Optical Test of Local Realistic Theories and Bell's Theorem}
It is possible to produce any of the four Bell states using either type I or type II down-conversion. For a birefringent crystal, the ordinary beam is along the path of the input beam while the extraordinary beam changes direction. The refraction is different for different polarizations. For photons passing through the birefringent crystal with an angle $\theta$ the o-ray will be in the state $\ket{\theta}$ while the e-ray will be in the state $\ket{\theta_\perp}$.

For two crystals with angle $\theta$ and $\phi$ that are placed in front of the ordinary or extraordinary beams, we can define the Bell state
\begin{align}
    \ket{\Psi^-} = \frac{1}{\sqrt{2}}  [ &(\cos\theta\sin\phi - \sin\theta\cos\phi) \ket{\theta}_1\ket{\phi}_2 + (\cos\theta\cos\phi + \sin\theta\sin\phi)\ket{\theta}_1\ket{\phi_\perp}_2\\
    -&(\sin\theta\sin\phi + \cos\theta\cos\phi)\ket{\theta_\perp}_1\ket{\phi}_2 - (\sin\theta\cos\phi - \cos\theta\sin\phi)\ket{\theta_\perp}_1\ket{\phi_\perp}_2].
\end{align}
For $\theta=\phi$ the Bell state $\ket{\Psi^-} = (\ket{\theta}_1\ket{\theta_\perp}_2 - \ket{\theta_\perp}_1\ket{\theta}_2)/\sqrt{2}$. We can test quantum mechanics against local hidden variable theories assuming space-light, space-time separation. Local hidden variable theories essentially say that we lack some information about the system (a hidden variable) that determines if the state is $\ket{\theta}_1\ket{\theta_\perp}_2$ or $\ket{\theta_\perp}_1\ket{\theta}_2$, and the measurement merely shows the outcome of this hidden variable. Another interpretation is that the act of measuring on of the photons forces the collapse of the wave function, and that this information somehow is transferred to the other photon, seemingly violating the speed of light, and is thus non-local.

We define two operators $\hat{A}(\theta)$ and $\hat{B}(\phi)$, with value $+1$ if the photon emerges in the ordinary beam and $-1$ if the photon emerges in the extraordinary beam for the two birefringent crystals, where $\theta,\phi$ is the angles for the different crystals. We then define the correlation function $C(\theta,\phi) = \langle \hat{A}(\theta)\hat{B}(\phi) \rangle$, which for the state $\ket{\Psi^-}$ is $C(\theta,\phi) = -\cos[2(\theta-\phi)]$, which is a quantum mechanical result. For a local hidden-variable theory we introduce the hidden variable $\lambda$ which affects the outcome of the measurement, and thus we have $\hat{A}(\theta,\lambda)$ and $\hat{B}(\phi,\lambda)$, with the values $\pm 1$ still. If $\density(\lambda)$ is the probability distribution of the hidden variable we obtain the correlation function in local hidden variable theory as $C_\text{HV}(\theta,\phi) = \int d\lambda \density(\lambda) \hat{A}(\theta,\lambda)\hat{B}(\phi,\lambda)$. Introducing the quantity 
\begin{equation}
    S = \int d\lambda\density(\lambda) (\hat{A}(\theta,\lambda) [\hat{B}(\phi,\lambda) + \hat{B}(\phi',\lambda)] +\hat{A}(\theta',\lambda) [\hat{B}(\phi,\lambda) - \hat{B}(\phi',\lambda)]) = \pm 2
\end{equation}
we have that
\begin{equation}
    -2\leq C_\text{HV}(\theta,\phi) + C_\text{HV}(\theta',\phi) + C_\text{HV}(\theta,\phi') - C_\text{HV}(\theta',\phi') \leq 2,
\end{equation}
which is the Clauser-Horne-Shimony-Holy-inequality (CHSH-inequality). However, choosing the angles $\theta=0, \theta'=\pi/4,\phi=\pi/8,$ and $\phi'=-\pi/8$ gives the quantum mechanical $S_\text{QM} = C(\theta,\phi) + C(\theta',\phi) + C(\theta,\phi') - C(\theta',\phi') = 2\sqrt{2}$, that is, quantum mechanics have results that no local hidden variable theory can explain.

\subsection{Franson's Experiment}
Imagine two photons, signal and idler, simultaneously being created by down-conversion and aimed in two directions. There is no single photon interference since the photon coherence times are much shorter than the difference in path lengths in the arms of the interferometer. The path difference in the arms comes from a beam splitter leading the photon to an adjustable length $L_{s/i}$ or a shorter length $S$. The paths ($L_s,L_i$) and $(S,S)$ are indistinguishable since the detectors click simultaneously, and we do not know when the photons were created, while $(L_s,S)$ and $(S,L_i)$ are distinguishable and rejected to obtain the reduced state $\ket{\psi} = (\ket{S}_s\ket{S}_i + e\exp(i\Phi)\ket{L_s}_s\ket{L_i}_i)/2$. The angle $\Phi$ is the relative phase between the two path processes, $\Phi \approx \omega_p (\Delta L_s +\Delta L_i)/2c$. We find that the probability of coincidence detection is $P_\text{coincidence} = [1 + \cos(\omega_p (\Delta L_s +\Delta L_i)/2c)]/2$. If the correlation fringes have a greater visibility than $\qty{70.7}{\percent}$ it violates Bell's inequality. Experiments have measures a visibility of $\qty{80.4(6)}{\percent}$. \question{I have a hard time completely understanding this experiment. But the results seem significant.}

\subsection{Application of Down-Converted Light to Metrology without Absolute Standards}
Down-conversion may be used to calculate the quantum efficiency of detectors without the use of a calibrated standard. NIST have tested this and found good agreement to more conventional method. \study{Since we essentially can use down-conversion to increase accuracy without a calibration, can we maybe use this in gate-operations in quantum computing as well.}